\documentclass[12pt,a4paper]{ctexart}
\usepackage{amsmath,amssymb}
\usepackage{geometry}
\geometry{left=2.5cm,right=2.5cm,top=2.5cm,bottom=2.5cm}
\usepackage{enumitem}
\usepackage{array}
\usepackage{booktabs}
\usepackage{fancyhdr}

% 设置页眉页脚
\pagestyle{fancy}
\fancyhf{} % 清除所有页眉页脚
\fancyfoot[C]{\thepage} % 页码居中放在页脚
\renewcommand{\headrulewidth}{0pt} % 去掉页眉横线

\title{回归系数的假设检验与区间估计——细致讲解}
\author{}
\date{}

\begin{document}

\maketitle

\section*{第一步：我们为什么需要“假设检验”和“区间估计”？}

想象一个场景：你想研究“家庭月收入”对“家庭月消费”的影响。你收集了10个家庭的数据，用最小二乘法（OLS）算出了一条回归直线：
\[
\text{消费} = 266.38 + 0.604 \times \text{收入}
\]
这里 \(0.604\) 就是斜率系数 \(\hat{\beta}_2\) 的估计值。它意味着：收入每增加1元，消费平均增加约0.604元。

但问题来了：
\begin{enumerate}[label=\arabic*.]
    \item 这个0.604是\textbf{真实}的影响吗？还是因为\textbf{抽样误差}偶然得到的？（比如换另一组10个家庭，算出来可能是0.58或0.62）
    \item 真实的 \(\beta_2\) 会不会其实是 \(0\)？如果真实是0，那就意味着\textbf{收入对消费根本没影响}，我们的0.604只是抽样的假象。
\end{enumerate}

\textbf{假设检验}就是为了回答第一个问题：这个估计值是否足够“显著”地不为零（或不为某个特定值）。\\
\textbf{区间估计}则是为了回答：真实的 \(\beta_2\) 大概落在什么范围内（比如0.565到0.642之间）。

\section*{第二步：OLS估计量 \(\hat{\beta}_1\) 和 \(\hat{\beta}_2\) 的分布性质（这是所有推断的基础）}

\subsection*{2.1 为什么它们是随机变量？}
因为我们用的是\textbf{样本}数据，每次抽样得到的数据不一样，算出来的 \(\hat{\beta}_1\) 和 \(\hat{\beta}_2\) 就不一样。所以它们都是\textbf{随机变量}，有自己的概率分布。

\subsection*{2.2 它们服从什么分布？}
在经典线性回归假设下，我们假定误差项 \(u_i\) 服从正态分布。因为 \(Y_i = \beta_1 + \beta_2 X_i + u_i\)，所以 \(Y_i\) 也服从正态分布。之前我们已经证明过，\(\hat{\beta}_1\) 和 \(\hat{\beta}_2\) 都是 \(Y_i\) 的线性组合（线性函数），而正态随机变量的线性组合仍然服从正态分布。

因此，即使样本很小，\(\hat{\beta}_1\) 和 \(\hat{\beta}_2\) \textbf{精确服从正态分布}。

（如果样本很大，即使误差项不是正态，根据中心极限定理，\(\hat{\beta}\) 也会\textbf{趋近}正态分布。）

\subsection*{2.3 正态分布的参数是什么？}
一个正态分布由两个参数决定：均值（期望）和方差。

我们已经证明过：
\begin{itemize}
    \item \(\hat{\beta}_1\) 和 \(\hat{\beta}_2\) 是\textbf{无偏估计量}，即 \(E(\hat{\beta}_1) = \beta_1\)，\(E(\hat{\beta}_2) = \beta_2\)。所以分布的\textbf{中心}就是真实参数值。
    \item 它们的方差公式也已经推导过（这里不重复推导，直接用结论）：
    \[
    \text{Var}(\hat{\beta}_1) = \sigma^2 \cdot \frac{\sum X_i^2}{n \sum x_i^2}
    \]
    \[
    \text{Var}(\hat{\beta}_2) = \frac{\sigma^2}{\sum x_i^2}
    \]
    其中 \(\sum x_i^2 = \sum (X_i - \bar{X})^2\)，\(\sigma^2\) 是误差项 \(u_i\) 的方差。
\end{itemize}

于是，它们的分布可写成：
\[
\hat{\beta}_1 \sim N\left(\beta_1, \sigma^2 \cdot \frac{\sum X_i^2}{n \sum x_i^2}\right)
\]
\[
\hat{\beta}_2 \sim N\left(\beta_2, \frac{\sigma^2}{\sum x_i^2}\right)
\]

\textbf{关键点}：这两个分布的中心是真实参数，离散程度由误差方差 \(\sigma^2\) 和样本数据决定。

\section*{第三步：标准化变换——从正态到标准正态}

为了做统计推断，我们通常把一般的正态变量转换成\textbf{标准正态变量} \(z\)，方法是：
\[
z = \frac{\text{随机变量} - \text{其期望}}{\text{其标准差}}
\]
这个变换后的 \(z\) 服从标准正态分布 \(N(0,1)\)（均值为0，方差为1）。

对 \(\hat{\beta}_2\) 做这个变换：
\begin{itemize}
    \item 随机变量：\(\hat{\beta}_2\)
    \item 期望：\(\beta_2\)
    \item 标准差：\(\text{SE}(\hat{\beta}_2) = \sqrt{\text{Var}(\hat{\beta}_2)} = \frac{\sigma}{\sqrt{\sum x_i^2}}\)
\end{itemize}

于是得到：
\[
z = \frac{\hat{\beta}_2 - \beta_2}{\text{SE}(\hat{\beta}_2)} = \frac{\hat{\beta}_2 - \beta_2}{\sigma / \sqrt{\sum x_i^2}} \sim N(0,1)
\]

同理，对 \(\hat{\beta}_1\) 也有类似结果。

\section*{第四步：现实中的麻烦——\(\sigma^2\) 未知怎么办？}

注意，上面那个 \(z\) 公式的分母里有 \(\sigma\)，它是\textbf{总体}误差的方差。现实中我们根本不知道 \(\sigma\) 的真实值，只能从样本去\textbf{估计}它。

我们用残差 \(e_i = Y_i - \hat{Y}_i\) 来估计误差，得到 \(\sigma^2\) 的无偏估计量：
\[
\hat{\sigma}^2 = \frac{\sum e_i^2}{n-2}
\]
（为什么除以 \(n-2\)？因为估计两个参数 \(\beta_1,\beta_2\) 用掉了两个自由度。）

然后，我们用这个估计值 \(\hat{\sigma}^2\) 来代替真实的 \(\sigma^2\)，计算出\textbf{估计的标准误差}（注意名字变化了）：
\[
\widehat{\text{SE}}(\hat{\beta}_2) = \frac{\hat{\sigma}}{\sqrt{\sum x_i^2}}
\]

这时候，如果我们还用这个估计的标准误差去做标准化：
\[
t = \frac{\hat{\beta}_2 - \beta_2}{\widehat{\text{SE}}(\hat{\beta}_2)}
\]
这个比值还服从标准正态分布吗？

\textbf{答案是：不服从了。} 因为分母不再是一个常数，而是一个\textbf{随机变量}（\(\hat{\sigma}\) 是随样本变化的）。

\subsection*{4.1 为什么变成了 \(t\) 分布？}
严格证明需要数学推导，但直观理解是：
\begin{itemize}
    \item 分子 \(\hat{\beta}_2 - \beta_2\) 服从正态分布。
    \item 分母涉及 \(\hat{\sigma}^2\)，而 \(\hat{\sigma}^2\) 与卡方分布有关。
    \item 一个正态随机变量除以一个独立的卡方随机变量（除以其自由度）的平方根，得到的比值服从 \textbf{学生氏 \(t\) 分布}。
\end{itemize}

在回归模型中，\(\hat{\beta}_2\) 与 \(\hat{\sigma}^2\) 是独立的，且 \(\hat{\sigma}^2\) 的自由度是 \(n-2\)，所以：
\[
t = \frac{\hat{\beta}_2 - \beta_2}{\widehat{\text{SE}}(\hat{\beta}_2)} \sim t(n-2)
\]
即自由度为 \(n-2\) 的 \(t\) 分布。

\textbf{注意}：
\begin{itemize}
    \item 如果样本量很大（\(n \to \infty\)），\(t\) 分布会趋近于标准正态分布。此时用 \(z\) 或 \(t\) 差别不大。
    \item 如果样本量小，必须用 \(t\) 分布，因为它的尾巴比正态分布厚，更保守，能反映用估计方差带来的额外不确定性。
\end{itemize}

\section*{第五步：回归系数的假设检验（核心实操）}

\subsection*{5.1 检验的基本思路}
我们想检验一个原假设（比如 \(\beta_2 = 0\)，即“X 对 Y 没有影响”）。如果这个假设是真的，那么从数据算出的 \(\hat{\beta}_2\) 应该离0不远；如果离0很远，我们就怀疑原假设不成立。

具体做法：
\begin{enumerate}
    \item 假设 \(H_0: \beta_2 = 0\) 为真。
    \item 计算检验统计量：\(t = \frac{\hat{\beta}_2 - 0}{\widehat{\text{SE}}(\hat{\beta}_2)}\)。
    \item 看这个 \(t\) 值落在 \(t\) 分布的哪个位置。如果落在两端极端区域（小概率事件发生），就拒绝原假设。
\end{enumerate}

\subsection*{5.2 双侧检验的步骤（以检验 \(\beta_2 = 0\) 为例）}
\textbf{第一步}：提出假设
\begin{itemize}
    \item 原假设 \(H_0: \beta_2 = 0\)
    \item 备择假设 \(H_1: \beta_2 \neq 0\)
\end{itemize}

\textbf{第二步}：计算 \(t\) 统计量
\begin{itemize}
    \item 用样本数据算出 \(\hat{\beta}_2\) 和它的标准误差 \(\widehat{\text{SE}}(\hat{\beta}_2)\)。
    \item 代入公式 \(t = \frac{\hat{\beta}_2}{\widehat{\text{SE}}(\hat{\beta}_2)}\)。
\end{itemize}

\textbf{第三步}：确定显著性水平 \(\alpha\)（通常取0.05或0.01）
\begin{itemize}
    \item \(\alpha\) 是我们愿意承担的“错误拒绝原假设”的最大风险。
    \item 对于双侧检验，我们在 \(t\) 分布左右两边各划出 \(\alpha/2\) 的拒绝域。
\end{itemize}

\textbf{第四步}：查表得到临界值
\begin{itemize}
    \item 自由度 \(df = n-2\)。
    \item 比如 \(n=10\)，\(df=8\)，\(\alpha=0.05\)，查 \(t\) 分布表得 \(t_{0.025}(8) = 2.306\)。
    \item 拒绝域是：\(t < -2.306\) 或 \(t > 2.306\)。
\end{itemize}

\textbf{第五步}：比较决策
\begin{itemize}
    \item 如果样本计算的 \(|t| > 2.306\)，则落在拒绝域，拒绝 \(H_0\)，认为 \(\beta_2\) 显著不为0。
    \item 如果 \(|t| \le 2.306\)，则不拒绝 \(H_0\)，认为没有足够证据说 \(\beta_2\) 不为0。
\end{itemize}

\subsection*{5.3 例子演练（消费-收入数据）}
已知：\(\hat{\beta}_2 = 0.60366\)，\(\widehat{\text{SE}}(\hat{\beta}_2) = 0.01672\)，\(n=10\)。

计算 \(t\) 值：
\[
t = \frac{0.60366}{0.01672} = 36.1085
\]
自由度 \(df = 10 - 2 = 8\)，临界值 \(t_{0.025}(8) = 2.306\)。

显然 \(36.1085\) 远大于 \(2.306\)，所以强烈拒绝 \(H_0: \beta_2 = 0\)，认为收入对消费有显著影响。

\subsection*{5.4 显著性水平变化会导致结论变化吗？}
例如：如果 \(t=2.5000\)，\(df=18\)。
\begin{itemize}
    \item 当 \(\alpha=0.05\)，临界值 \(2.1009\)，\(|t|>2.1009\)，拒绝 \(H_0\)。
    \item 当 \(\alpha=0.01\)，临界值 \(2.8784\)，\(|t|<2.8784\)，不拒绝 \(H_0\)。
\end{itemize}
这说明：\textbf{结论依赖于你选的严格程度 \(\alpha\)}。\(\alpha\) 越小，越不容易拒绝原假设。

\subsection*{5.5 更方便的方法：\(P\) 值检验}
为了避免每次查表，软件通常给出 \(P\) 值。

\textbf{定义}：在 \(H_0\) 为真的条件下，观察到当前样本 \(t\) 统计量绝对值（或更极端情况）的概率。
\begin{itemize}
    \item 对于双侧检验，\(P\) 值 = \(P(|T| > |t_{\text{样本}}|)\)。
\end{itemize}

\textbf{决策规则}：
\begin{itemize}
    \item 如果 \(P\) 值 < \(\alpha\)，拒绝 \(H_0\)。
    \item 如果 \(P\) 值 > \(\alpha\)，不拒绝 \(H_0\)。
\end{itemize}

例如，前面 \(t=36.1085\) 对应的 \(P\) 值几乎是0.0000，远小于0.05，所以拒绝 \(H_0\)。

\section*{第六步：回归系数的区间估计（置信区间）}

\subsection*{6.1 区间估计想干什么？}
点估计（如 \(\hat{\beta}_2=0.604\)）只是单个数值，我们不知道它离真实值有多近。区间估计则是给出一个\textbf{范围}，并且说明这个范围包含真实值的概率（置信水平）。

公式形式：
\[
\hat{\beta}_2 \pm \text{临界值} \times \widehat{\text{SE}}(\hat{\beta}_2)
\]

\subsection*{6.2 小样本下（\(\sigma^2\) 未知）的置信区间推导}
因为 \(t = \frac{\hat{\beta}_2 - \beta_2}{\widehat{\text{SE}}(\hat{\beta}_2)} \sim t(n-2)\)，

我们可以写出概率表达式：
\[
P\left( -t_{\alpha/2} < \frac{\hat{\beta}_2 - \beta_2}{\widehat{\text{SE}}(\hat{\beta}_2)} < t_{\alpha/2} \right) = 1 - \alpha
\]
其中 \(t_{\alpha/2}\) 是自由度为 \(n-2\) 的 \(t\) 分布双侧临界值。

把不等式变形，解出 \(\beta_2\)：
\begin{enumerate}
    \item 乘以标准误差（正数，不等号方向不变）：
    \[
    -t_{\alpha/2} \cdot \widehat{\text{SE}}(\hat{\beta}_2) < \hat{\beta}_2 - \beta_2 < t_{\alpha/2} \cdot \widehat{\text{SE}}(\hat{\beta}_2)
    \]
    \item 各项减去 \(\hat{\beta}_2\)：
    \[
    -\hat{\beta}_2 - t_{\alpha/2} \cdot \widehat{\text{SE}}(\hat{\beta}_2) < -\beta_2 < -\hat{\beta}_2 + t_{\alpha/2} \cdot \widehat{\text{SE}}(\hat{\beta}_2)
    \]
    \item 乘以 \(-1\)（不等号反向）：
    \[
    \hat{\beta}_2 + t_{\alpha/2} \cdot \widehat{\text{SE}}(\hat{\beta}_2) > \beta_2 > \hat{\beta}_2 - t_{\alpha/2} \cdot \widehat{\text{SE}}(\hat{\beta}_2)
    \]
    \item 重新排列，得到熟悉的置信区间：
    \[
    \hat{\beta}_2 - t_{\alpha/2} \cdot \widehat{\text{SE}}(\hat{\beta}_2) < \beta_2 < \hat{\beta}_2 + t_{\alpha/2} \cdot \widehat{\text{SE}}(\hat{\beta}_2)
    \]
\end{enumerate}

因此，\textbf{置信水平为 \(1-\alpha\) 的置信区间是}：
\[
\left[ \hat{\beta}_2 - t_{\alpha/2}(n-2) \cdot \widehat{\text{SE}}(\hat{\beta}_2),\quad \hat{\beta}_2 + t_{\alpha/2}(n-2) \cdot \widehat{\text{SE}}(\hat{\beta}_2) \right]
\]

\subsection*{6.3 例子计算}
消费-收入数据：\(\hat{\beta}_2 = 0.60366\)，\(\widehat{\text{SE}} = 0.01672\)，\(n=10\)，取 \(\alpha=0.05\)。
\begin{itemize}
    \item \(df = 8\)，\(t_{0.025}(8) = 2.306\)。
    \item 误差幅度 = \(2.306 \times 0.01672 = 0.03855\)。
    \item 下限 = \(0.60366 - 0.03855 = 0.56511\)。
    \item 上限 = \(0.60366 + 0.03855 = 0.64221\)。
\end{itemize}
所以，\(\beta_2\) 的95\%置信区间为 \([0.5651, 0.6422]\)。

\subsection*{6.4 对置信区间的正确理解}
\begin{itemize}
    \item \textbf{不是}说“\(\beta_2\) 有95\%的概率落在0.5651到0.6422之间”。真实 \(\beta_2\) 是一个固定但未知的常数，它要么在区间里，要么不在，概率只能是0或1。
    \item \textbf{正确解释}：如果我们重复抽样100次，每次按同样方法计算一个95\%置信区间，那么大约有95个区间会包含真实的 \(\beta_2\)。我们手头的这个区间是这100次中的一次实现，我们有95\%的信心说它盖住了真实值。
\end{itemize}

\section*{总结对照表}

\begin{table}[h]
\centering
\caption{不同情形下的推断方法对比}
\begin{tabular}{>{\raggedright\arraybackslash}p{2.5cm}>{\raggedright\arraybackslash}p{2.5cm}>{\raggedright\arraybackslash}p{2cm}>{\raggedright\arraybackslash}p{3cm}>{\raggedright\arraybackslash}p{2.5cm}>{\raggedright\arraybackslash}p{2.5cm}}
\toprule
\textbf{情形} & \textbf{方差 \(\sigma^2\)} & \textbf{样本量} & \textbf{统计量及其分布} & \textbf{用于假设检验} & \textbf{用于区间估计} \\
\midrule
理想情况 & 已知 & 任意 & \(z \sim N(0,1)\) & 查正态表 & 正态分位数 \\
大样本 & 未知（用 \(\hat{\sigma}^2\) 估计） & 大 & 近似 \(z \sim N(0,1)\) & 查正态表 & 正态分位数 \\
小样本（最常见） & 未知（用 \(\hat{\sigma}^2\) 估计） & 小 & \(t \sim t(n-2)\) & 查 \(t\) 表 & \(t\) 分位数 \\
\bottomrule
\end{tabular}
\end{table}

通过以上极其细致的步骤，你应该能完全理解回归系数的假设检验和区间估计的来龙去脉，以及为什么在小样本中必须使用 \(t\) 分布。

\end{document}